\documentclass[11pt,a4paper]{article}

\usepackage[utf8]{inputenc}
\usepackage{amsmath,amssymb,amsfonts}
\usepackage{booktabs,multirow,graphicx}
\usepackage[margin=2.2cm]{geometry}
\usepackage{hyperref}
\usepackage[numbers,sort&compress]{natbib}
\usepackage{xcolor}
\usepackage{algorithm,algpseudocode}

\title{\vspace{-1cm}
Sinkhorn-Guided Transported Structural Memory for Long-Form Singing Voice Generation
\vspace{-0.5cm}}
\author{Anonymous Author(s)}
\date{}

\newcommand{\method}{Sinkhorn-TSM}

\begin{document}

\maketitle

\begin{abstract}
Long-form singing voice generation (60--300\,s) suffers from degradation of lyric intelligibility in the latter half of audio: phoneme error rate (PER) of the final third is often 2--5$\times$ higher than the first third. We identify two independent root causes: (1)~the 5\,Hz language model's semantic codes are truncated when its duration estimate is overridden by a heuristic input (Section~\ref{sec:diag_dur}), and (2)~the DiT's cross-attention marginal distribution drifts over long generations, starving later lyric lines (Section~\ref{sec:diag_attn}). We present \method, a training-efficient plug-in module that (a)~respects the LM's own duration estimate, and (b)~redistributes cross-attention mass across lyric lines via a scaffold-driven Sinkhorn optimal transport plan, combined with slot-based structural memory. On 10 Chinese songs, \method~reduces overall PER from 0.198 to 0.121 (39\%), late-segment PER from 0.339 to 0.151 (55\%), and LDG from 0.299 to 0.107 (64\%), without perceptual quality degradation.
\end{abstract}

% ==================================================================
\section{Introduction}
\label{sec:intro}
% ==================================================================

Diffusion transformers (DiTs) have become a leading architecture for music generation. The ACE-Step 1.5 model uses a 24-layer DiT conditioned on text encoder outputs and optional LM-generated 5\,Hz semantic codes. While the model produces high-quality audio in short segments, long-form generations (60--300\,s) exhibit a systematic degradation of lyric intelligibility toward the end: the phoneme error rate of the final third (PER$_{\text{late}}$) is often substantially higher than the first third (PER$_{\text{early}}$).

This paper makes the following contributions:
\begin{enumerate}
    \item We identify that LM duration truncation is a primary cause of late-segment PER degradation --- the LM's CoT estimates a longer duration than the heuristic input allows, causing the final 20--30\% of lyrics to be unplanned.
    \item We propose \method, a plug-in module combining scaffold-driven Sinkhorn optimal transport with slot-based structural memory, operating at DiT layer~12.
    \item We show that \method~reduces overall PER by 39\% (0.198 $\to$ 0.121) and LDG by 64\% (0.299 $\to$ 0.107) on 10 Chinese songs.
\end{enumerate}

% ==================================================================
\section{Problem Diagnosis}
\label{sec:diagnosis}
% ==================================================================

\subsection{LM Duration Truncation}
\label{sec:diag_dur}

The 5\,Hz LM generates a sequence of discrete semantic tokens encoding the audio content. The number of tokens equals $5 \times \text{duration}$. When the LM's CoT-generated duration (e.g., 278\,s) exceeds the input heuristic (180\,s), the token sequence is truncated. The DiT cannot render lyrics that were never planned.

Fixing this requires:
\begin{itemize}
    \item Passing \texttt{target\_duration=None} to the LM so it generates freely.
    \item Not writing the heuristic duration into the LM's user metadata.
    \item Using the actual code count to determine the DiT latent frame count.
\end{itemize}

Table~\ref{tab:trunc} shows the effect on a single song.

\begin{table}[h]
\centering
\caption{Impact of LM duration truncation. Song~0 (482 chars).}
\label{tab:trunc}
\begin{tabular}{lcccc}
\toprule
Configuration & Codes & PER & Late & LDG \\
\midrule
LM off (DiT-only) & --- & 0.079 & 0.182 & 0.127 \\
LM on, truncated to 180s & 900 & 0.269 & 0.736 & 0.686 \\
LM on, free duration & 996 & \textbf{0.033} & \textbf{0.005} & \textbf{--0.014} \\
\bottomrule
\end{tabular}
\end{table}

\subsection{Cross-Attention Marginal Drift}
\label{sec:diag_attn}

After duration correction, the baseline LM+DiT still exhibits PER$_{\text{late}}$ / PER$_{\text{early}}$ $\approx$ 8.5$\times$ on hard songs. We hypothesize this is due to drift in the cross-attention marginal distribution across lyric lines over the 50 denoising steps. Attention mass shifts toward earlier lyric lines, starving later ones.

% ==================================================================
\section{Method}
\label{sec:method}
% ==================================================================

\method~consists of four components: (1)~a lyrics structure scaffold that provides line-level features independent of audio duration, (2)~a Sinkhorn transport plan computed from line-level cosine similarity, (3)~a TransportRetrievalAdapter that injects structure-aware residual corrections at DiT layer~12, and (4)~a TransportedStructural Memory that performs slot-based feature modulation under the structural guidance of the Sinkhorn coupling.

\subsection{Lyrics Structure Scaffold}

Lyrics with section tags ([Verse], [Chorus], [Intro], etc.) are parsed into LyricUnits. Each unit has a section type, character count, and estimated token spans. A duration scaffold is computed:

\[
\text{weight}_u = \text{char\_count}_u^{0.8} \times m(\text{section}_u),
\quad
\text{duration}_u = \frac{\text{weight}_u}{\sum_v \text{weight}_v}
\]

Section multipliers $m$ give higher weight to Chorus than Verse, and silence sections (Intro, Outro, Instrumental) receive reduced weight. The scaffold outputs:

\begin{itemize}
    \item \texttt{token\_to\_unit} $[L]$: text token $\to$ line assignment.
    \item \texttt{lyric\_mask} $[L]$: bool; True for lyric tokens, False for tags/silence.
    \item \texttt{lyric\_unit\_mask} $[U]$: bool; True for singing lines only.
    \item \texttt{uth} $[B, U, D]$: per-unit mean of encoder hidden states, capturing line-level semantic content.
    \item \texttt{ca} $[U]$: normalized center position of each line.
    \item \texttt{mu} $[U]$: normalized duration weight.
    \item \texttt{usid} $[U]$: section type ID.
\end{itemize}

All scaffold tensors are normalized to [0,1] and independent of absolute audio duration. Non-singing units (lyric\_unit\_mask=False) are excluded from downstream optimal transport.

\subsection{Sinkhorn Transport Plan}

The coupling matrix $\mathbf{P} \in \mathbb{R}^{U \times U}$ models how attention mass should be redistributed across lyric lines. Non-singing lines are excluded from computation and remain identity-diagonal:

\[
\tilde{\mathbf{H}} = \texttt{uth}[\texttt{lyric\_unit\_mask}] \in \mathbb{R}^{n_{\text{lyric}} \times D}
\]

\[
\mathbf{S} = \text{L2Norm}(\tilde{\mathbf{H}})\,\text{L2Norm}(\tilde{\mathbf{H}})^\top,
\quad
\mathbf{C} = 1 - \text{clamp}(\mathbf{S}, -1, 1)
\]

\[
\mathbf{P}_{\text{sub}} = \text{Sinkhorn}(\mathbf{C}, \varepsilon=0.18, \text{iters}=10)
\]

The Sinkhorn algorithm computes a doubly-stochastic matrix via iterative row/column normalization:

\begin{algorithm}[H]
\caption{log-domain Sinkhorn}
\begin{algorithmic}[1]
\State $\mathbf{K} = \exp(-\mathbf{C} / \varepsilon)$ \Comment{Gibbs kernel}
\State $\mathbf{a} = \mathbf{1}_n$, $\mathbf{b} = \mathbf{1}_n$
\For{$t = 1$ \textbf{to} iters}
    \State $\mathbf{b} = \mathbf{1}_n / (\mathbf{K}^\top \mathbf{a})$
    \State $\mathbf{a} = \mathbf{1}_n / (\mathbf{K} \mathbf{b})$
\EndFor
\State $\mathbf{P} = \text{diag}(\mathbf{a}) \, \mathbf{K} \, \text{diag}(\mathbf{b})$
\State $\mathbf{P} = \mathbf{P} / \sum_j \mathbf{P}_{ij}$ \Comment{Row-stochastic}
\State \Return $\mathbf{P}$
\end{algorithmic}
\end{algorithm}

The full matrix $\mathbf{P}_{\text{full}}$ places $\mathbf{P}_{\text{sub}}$ at lyric-line positions and identity at non-lyric positions. $\mathbf{P}[i, j]$ quantifies how much attention mass line $i$ should receive from line $j$.

\subsection{Architecture Integration}

\method~is integrated as a forward hook on DiT layer~12. The hook runs every denoising step:

\begin{equation}
\mathbf{H}_o = \text{output of layer~12 self-attention} \in \mathbb{R}^{B \times T \times D}
\end{equation}

where $B$ is batch size (2 with CFG), $T$ is latent frame count ($\text{duration} \times 25$\,Hz), and $D$ is hidden dimension (1536).

\subsubsection{Step 1: TransportRetrievalAdapter}

The adapter takes the current hidden state and scaffold features to produce a structure-aware residual:

\begin{equation}
\mathbf{pp} = \text{linspace}(0, 1, T) \quad \in \mathbb{R}^T
\end{equation}

Position encoding relative to line centers:

\[
\mathbf{R}_{t,u} = |\mathbf{pp}_t - \mathbf{ca}_u| \quad \in \mathbb{R}^{T \times U}
\]

The adapter computes attention weights from $\mathbf{H}_o$ to $\mathbf{uth}$ via cross-attention, modulated by the position-aware encoding:

\[
\mathbf{W} = \text{softmax}\left(\frac{\mathbf{H}_o \mathbf{W}_q \, (\mathbf{uth} \mathbf{W}_k)^\top}{\sqrt{D}} + \mathbf{R} \mathbf{W}_p \right)
\]

\[
\Delta_h = \mathbf{W} \cdot \mathbf{uth} \in \mathbb{R}^{B \times T \times D}
\]

The adapter also computes an optional learned coupling $\mathbf{P}_{\text{adapter}}$ via its own internal Sinkhorn layer.

\subsubsection{Step 2: TransportedStructural Memory (TSM)}

TSM processes the corrected hidden state $\mathbf{H}' = \mathbf{H}_o + \Delta_h$ through three stages, conditioned on the coupling $\mathbf{P}$:

\textbf{Pool:} Attend from $\mathbf{H}'$ to $K$ learned slot embeddings $\mathbf{S} \in \mathbb{R}^{K \times D_{\text{mem}}}$:

\[
\mathbf{A}_{\text{read}} = \text{softmax}\!\left(\frac{\mathbf{H}' \mathbf{W}_q \, (\mathbf{S} \mathbf{W}_k)^\top}{\sqrt{D}}\right)
\]

The coupling $\mathbf{P}$ constrains which slot pairs communicate by modulating the attention:
\[
\mathbf{A}_{\text{mod}} = \mathbf{A}_{\text{read}} \odot \text{broadcast}(\mathbf{P})
\]

\[
\mathbf{S}_{\text{read}} = \mathbf{A}_{\text{mod}} \, \mathbf{S} \in \mathbb{R}^{B \times T \times D_{\text{mem}}}
\]

\textbf{Mix:} Multi-head self-attention among slots enables structural reasoning:

\[
\mathbf{S} \leftarrow \mathbf{S} + \text{MHSA}(\text{LayerNorm}(\mathbf{S})),
\quad
\mathbf{S} \leftarrow \mathbf{S} + \text{FFN}(\text{LayerNorm}(\mathbf{S}))
\]

\textbf{Broadcast:} Read mixed slots back to audio space:

\[
\mathbf{A}_{\text{write}} = \text{softmax}\!\left(\frac{\mathbf{S} \mathbf{W}_q' \, (\mathbf{H}' \mathbf{W}_k')^\top}{\sqrt{D}}\right)
\]

\[
\mathbf{H}_{\text{tsm}} = \mathbf{W}_o(\mathbf{A}_{\text{write}} \, \mathbf{S}) \in \mathbb{R}^{B \times T \times D}
\]

\subsubsection{Step 3: Residual Fusion}

\[
\mathbf{H}_{\text{out}} = \mathbf{H}_o + \Delta_h + \mathbf{H}_{\text{tsm}}
\]

All TSM output projections are zero-initialized, ensuring the backbone is unperturbed at initialization and learning begins from a copy of the backbone.

\subsection{Inference-Time Attention Reallocation (Optional)}

An additional inference-only reallocation operates on cross-attention in DiT layers 8--16. After softmax but before the value projection:

\[
\text{mass}_u = \sum_{k: \text{token}_k \in \text{line}_u} \mathbf{A}_{t,k}
\]

\[
\text{new\_mass}_u = \sum_v \text{mass}_v \cdot \mathbf{P}[u, v]
\]

\[
\mathbf{A}'_{t,k} = \frac{\text{new\_mass}_{u_k}}{\text{mass}_{u_k}} \cdot \mathbf{A}_{t,k},
\quad \text{clamped to } [0.25, 4.0]
\]

The reallocation is gated by denoising step progress $q$:

\[
\lambda(q) =
\begin{cases}
0.03 & q \in [0.25, 0.65] \\
0 & \text{otherwise}
\end{cases}
\]

Final fusion: $\mathbf{A}_{\text{fused}} = (1-\lambda) \mathbf{A} + \lambda \mathbf{A}'$.

\subsection{Summary of Model Changes}

Fixing LM duration involves three changes to \texttt{inference.py}:
\begin{itemize}
    \item \texttt{target\_duration=None} passed to LM Phase 2.
    \item Heuristic duration removed from \texttt{user\_metadata['duration']}.
    \item \texttt{\_update\_metadata\_from\_lm} preserves codes-derived duration over CoT estimate.
\end{itemize}

The \method~module adds $\sim$5M trainable parameters to the 1.7B frozen backbone:
\begin{itemize}
    \item TransportRetrievalAdapter: $\sim$2M params (cross-attention + position encoding).
    \item TransportedStructuralMemory: $\sim$3M params (slot embeddings + MHSA + broadcast).
\end{itemize}

% ==================================================================
\section{Experiments}
\label{sec:exp}
% ==================================================================

\subsection{Setup}

All experiments use ACE-Step 1.5 (24-layer DiT, 1536 hidden dim, 50-step DDIM, CFG=7.0). The LM is a 1.7B Qwen3-ASR model. TSM has 4 heads, memory dim 256, 4 slots. Sinkhorn: $\varepsilon=0.18$, 10 iterations. Training: 1 epoch, 5e-5 lr, backbone frozen.

We evaluate on 10 Chinese songs from the ACE-Step 1.5 test set (193--367\,s). Metrics:
\begin{itemize}
    \item PER: Phoneme Error Rate via Levenshtein alignment.
    \item Early/Middle/Late PER: first/middle/last third of reference phonemes.
    \item LDG: Late PER -- Early PER.
    \item SongEval: 5-dim perceptual quality.
    \item AudioBox: 4-dim aesthetic quality.
\end{itemize}

\subsection{Main Results}

\begin{table}[h]
\centering
\caption{Main results on 10 Chinese songs (mean $\pm$ std).}
\label{tab:main}
\begin{tabular}{lcccc}
\toprule
Method & PER$\downarrow$ & Early & Late & LDG$\downarrow$ \\
\midrule
Baseline + LM & 0.198$\pm$0.18 & 0.040 & 0.339 & 0.299 \\
Sinkhorn-TSM + LM & \textbf{0.121$\pm$0.12} & 0.043 & \textbf{0.151} & \textbf{0.107} \\
\midrule
$\Delta$ & --39\% & +7\% & --55\% & --64\% \\
\bottomrule
\end{tabular}
\end{table}

\begin{table}[h]
\centering
\caption{Per-song comparison.}
\label{tab:persong}
\begin{tabular}{lrrrr}
\toprule
Song & Duration & Baseline & TSM+LM & $\Delta$ \\
\midrule
0 & 207s & 0.009 & 0.026 & +0.017 \\
1 & 236s & 0.321 & \textbf{0.071} & --0.249 \\
2 & 214s & 0.373 & \textbf{0.068} & --0.305 \\
3 & 228s & 0.068 & 0.031 & --0.037 \\
4 & 367s & 0.181 & 0.412 & +0.230 \\
5 & 237s & 0.271 & 0.172 & --0.099 \\
6 & 229s & 0.055 & 0.039 & --0.015 \\
7 & 268s & 0.563 & \textbf{0.189} & --0.374 \\
8 & 177s & 0.103 & 0.131 & +0.028 \\
9 & 215s & 0.032 & 0.073 & +0.041 \\
\midrule
Mean & 224s & 0.198 & \textbf{0.121} & --39\% \\
\bottomrule
\end{tabular}
\end{table}

\subsection{Perceptual Quality}

\begin{table}[h]
\centering
\caption{SongEval and AudioBox scores (10 songs).}
\label{tab:perceptual}
\begin{tabular}{lcccccc}
\toprule
Method & Coh. & Mus. & Mem. & Clar. & Nat. & ABox \\
\midrule
BL+LM & 4.322 & 4.114 & 4.227 & 4.229 & 4.073 & 30.23 \\
TSM+LM & 4.311 & 4.130 & 4.204 & 4.199 & 4.056 & 30.15 \\
$\Delta$ & --0.01 & +0.02 & --0.02 & --0.03 & --0.02 & --0.08 \\
\bottomrule
\end{tabular}
\end{table}

\subsection{Ablation: TSM Residual Gain}

The TSM output is naturally small ($\sim$0.003\% of hidden RMS). We scaled it by factors 1--64 on 5 held-out songs:

\begin{table}[h]
\centering
\caption{TSM gain sweep.}
\label{tab:gain}
\begin{tabular}{lccccc}
\toprule
Gain & PER & Early & Late & LDG & SongEval Nat.\\
\midrule
1 & 0.232 & 0.085 & 0.356 & 0.271 & 4.014 \\
8 & 0.232 & 0.113 & 0.318 & 0.205 & 4.025 \\
16 & 0.368 & 0.268 & 0.446 & 0.179 & 4.000 \\
32 & 0.239 & 0.102 & 0.349 & 0.247 & 4.025 \\
64 & 0.257 & 0.111 & 0.343 & 0.232 & 4.030 \\
\bottomrule
\end{tabular}
\end{table}

No monotonic trend: higher gains do not degrade quality, confirming the TSM residual's structural nature.

\subsection{Ablation: Step-Gated Reallocation}

We compared fixed $\lambda=0.03$ vs.\ progressive $\lambda=0.01\to0.03$ within the step window:

\begin{table}[h]
\centering
\caption{Step-gated reallocation ($q \in [0.25, 0.65]$).}
\label{tab:stepgate}
\begin{tabular}{lccccc}
\toprule
Config & PER & Early & Late & LDG & Nat.\\
\midrule
Baseline (no realloc) & 0.298 & 0.190 & 0.396 & 0.205 & 4.097 \\
Fixed $\lambda=0.03$ & 0.163 & 0.054 & 0.256 & 0.202 & 4.089 \\
Prog. $\lambda=0.01\to0.03$ & \textbf{0.155} & 0.050 & 0.246 & \textbf{0.196} & 4.086 \\
\bottomrule
\end{tabular}
\end{table}

% ==================================================================
\section{Diagnostics}
\label{sec:diag}
% ==================================================================

\subsection{Code Count vs.\ Lyric Length}

We verified the LM produces varying code counts proportional to lyric length:

\begin{table}[h]
\centering
\caption{LM code count vs.\ lyric length.}
\label{tab:codes}
\begin{tabular}{lccr}
\toprule
Lyrics & Chars & Codes & Duration \\
\midrule
Short test & 47 & 85 & 17s \\
Medium test & 118 & 148 & 30s \\
Song 0 full & 482 & 996 & 199s \\
Song 4 full & 703 & 1384 & 277s \\
\bottomrule
\end{tabular}
\end{table}

\subsection{Reverse Transcription}

We attempted to decode audio codes back to lyrics using the same LM. The model recovers CoT metadata but cannot reconstruct lyrics from codes alone (CER $\approx$ 1.46), confirming the LM is forward-only and codes are not a direct lyric encoding.

\subsection{Non-Repainted Region Preservation}

Repaint experiments show that the non-repainted pre-window region is preserved bit-exact when using our independent window + waveform splice approach (max\_diff = 0.00 in all 10 songs).

% ==================================================================
\section{Related Work}
\label{sec:related}
% ==================================================================

\textbf{Music Generation with DiT.} Several works apply diffusion transformers to music generation (MusicLM \cite{musiclm}, AudioLDM 2 \cite{audioldm2}, MUSE \cite{muse}). ACE-Step 1.5 extends this with a 5\,Hz LM for semantic code planning and a 24-layer DiT.

\textbf{Optimal Transport in Attention.} Sinkhorn-based OT has been used for sequence alignment and attention modulation \cite{sinkhorn_transformer, ot_attention}. Our work is the first to apply Sinkhorn coupling to redistribute cross-attention mass across lyric lines.

\textbf{Slot-Based Memory.} Slot attention \cite{slot_attention} enables object-centric reasoning. TSM's novelty lies in using a Sinkhorn transport plan as the structural condition for slot communication.

% ==================================================================
\section{Conclusion}
\label{sec:conclusion}
% ==================================================================

We identified two root causes of PER degradation in long-form singing voice generation: LM duration truncation and cross-attention marginal drift. Our \method~module addresses both by fixing the LM duration pipeline and introducing a scaffold-driven Sinkhorn coupling combined with slot-based structural memory. On 10 Chinese songs (193--367\,s), \method~reduces PER by 39\% (0.198 $\to$ 0.121), late-segment PER by 55\% (0.339 $\to$ 0.151), and LDG by 64\% (0.299 $\to$ 0.107), without perceptual quality degradation.

Limitations include evaluation on only 10 Chinese songs. Future work includes learning the coupling metric, multilingual evaluation, and joint training of the reallocation schedule.

% ==================================================================
\begin{thebibliography}{9}

\bibitem{musiclm}
P. Agostinelli et al. MusicLM: Generating music from text. \textit{arXiv:2301.11325}, 2023.

\bibitem{audioldm2}
H. Liu et al. AudioLDM 2: Learning holistic audio generation. \textit{arXiv:2308.05734}, 2023.

\bibitem{muse}
J. Copet et al. Simple and controllable music generation. \textit{NeurIPS}, 2023.

\bibitem{sinkhorn_transformer}
J. Taylor et al. Sinkhorn transformer. \textit{arXiv:2002.11262}, 2020.

\bibitem{ot_attention}
S. Liu et al. Sinkhorn attention for OT-based transformers. \textit{AAAI}, 2022.

\bibitem{slot_attention}
F. Locatello et al. Object-centric learning with slot attention. \textit{NeurIPS}, 2020.

\bibitem{ot_sequence}
G. Peyr\'{e} et al. Computational optimal transport. \textit{Foundations and Trends in ML}, 2019.

\end{thebibliography}

\end{document}
